\chapter{Conclusion and Future Work}
\label{ch:6}

\section{Summary of Findings}
\label{sec:6.1}

The goal of this research was to create an empathetic response system that is separated between emotion detection and response generation, to fit the entire system within the 8 GB of an ordinary consumer graphics card, and to measure whether or not the separation is worthwhile. From the 4,200 scored replies in fourteen configurations, the answer is a qualified yes, and the qualification is as straightforward as the support. The work is split in two parts. The first constructed the detector. Five encoders and decoders were trained in the same way, and the results of RoBERTa-base were chosen: 50.3\% top-1 accuracy and 84.3\% in the top-5 for the 32 emotions of EmpatheticDialogues, outperforming all jointly trained models published on this corpus. There are two findings that are more important than that number. Near neighbours are almost always the only mistakes the detector can make, for example, afraid is read as terrified, and sad is read as devastated; and in fact, what it is trained on is more important than which network is selected, as training on the situation descriptions of the corpus while using it on the user's own message cost around nine points. Both lead to passing three candidate emotions instead of one.

The second part constructed and tested the pipeline. Five different methods of connecting the two models were measured and compared to the detector, with the prompt fixed, and three generators were compared with the detector and prompt held fixed. The chosen design, which sees Phi-4-mini first select one of the three detector candidates and then write under per-emotion guidance, comes in second place of the eight systems with an empathy score of +0.59 versus +1.11 for a single model that is nearly double the size. The argument is carried in two comparisons. The same generator alone scores +0.53, so the gain is due to the design rather than the size of the generator; and the system scores 11.0 points higher than a single model of similar size, and 36.7\% of the replies are the speaker's true emotion. It does this using 3.40 GB of peak memory, 3.96 billion parameters and 4.17 seconds per reply, which is 38\% less memory than a single model, with 1.9 times the parameters and at no significant cost in response time. The human study confirms the comparison of the automatic measures that was left open: 40 participants rated the two-stage system highest on all three scales and preferred it in 61.3\% of the choices, compared to a single model that is almost twice the size in 21.6\% of the choices.

There are four findings that qualify the claim. The biggest single model is still leading on the empathy panel, and if it's deployed with the hardware to run a 7-billion-parameter model, and no need for an auditable emotional decision, it should use one. The system has its own three response rules, which takes up 60.0\% of the prompt space, versus 73.7\% for the same generator alone, as the rules share prompt space with the emotional guidance. Whether or not the two model judges distinguish the two configurations, the case is decided based on the empathy classifier, emotional fidelity, lower repetition and the visible decision, not on judge preference. But with the detector replaced by a more accurate one, three points better on the same test, no measure of reply quality changed, meaning that the detector was not the limiting component.

\section{Contributions to the Field}
\label{sec:6.2}

There are six contributions to this research. 
\textbf{First and foremost,} it gives you a calibrated understanding of what you're paying for and what you're paying for in a scale where it actually matters when you go for a "decoupled" design. The combination of a small emotion model and a large language model has been proposed earlier, but not tested against a hardware constraint of models of similar size and larger models, as this is the real-world choice one might face when deciding how to construct such a model.

\textbf{Second,} it makes note of negative outcomes, which can save other research workers time. A detector which was 3 points more accurate showed no improvement on any measure of the reply, suggesting the value of the detector is in the structure it provides, not in the precision of the label, and that further accuracy is not where the remaining gains lie. Similarly, retrieval augmentation (which was added on the reasonable expectation that related passages would help) led to a decrease in adherence to rules, from 63.7\% to 43.7\%, and did not help performance.

\textbf{Third,} it allows the emotional decision to be heard. Each response includes the three candidate emotions, the selected emotion and the phrases that triggered the emotion. An end-to-end model can provide a similar response, but it will not indicate its response -- and if a response is incorrect, there is no mechanism to secure a safeguard and no mechanism to review: in a sensitive domain, a difference of kind, not just degree.

\textbf{Fourth,} it provides an evaluation design that can be re-used in other work on this task. It does not rely on similarity to a single reference reply, but rather integrates three elements: explicit evaluation of the parts that a supportive reply must include; an empathy panel that compares two model judges from different developers and compares a classifier trained directly from human ratings; and paired statistical Testing with correction for multiple comparisons, which renders a difference as real only if it is detected by both the modified test and the confidence interval.

\textbf{Fifth}, it does not provide a benchmark entry; the context window for the detector, the memory for the generator and rules are relaxed after the first turn, preventing an interrogation of the conversation.

\textbf{Sixth,} it records the practical failures that a reproduction will reproduce, such as a fine-tuning defect: a reproduction of the model will not learn to stop if the wrong padding token is used, along with an attention implementation that does not return any weights at all.

\section{Limitations}
\label{sec:6.3}

There were seven limitations that restricted these conclusions as described in detail in Section 5.6.5. This is an English-language, crowd-sourced corpus with one label per conversation, so nothing said here should be taken as representative of Bangla language speakers or persons in actual distress. Label set doesn't include a category for suicidal ideation and there's no crisis safeguard in the system. All automatic measures are proxies: the acknowledgement check is looking for keywords, the empathy classifier is trained in a different domain and favours longer answers, emotional fidelity is being assessed by our own detector. The evidence is based on 300 conversations, 64 of which were viewed by the judges, and a small study of 40 people, whose ratings matched. All quantitative results are related to first replies as the multi-turn agent was not measured, but checked on scenarios. The figure for repetition is still high with 58.7\% of replies. The cost figures are from one graphic card, more than 50 messages.

\section{Recommendations for Future Work}
\label{sec:6.4}

\textbf{Immediate extensions:} A crisis filter that checks the message prior to generation and provides answers from a set of reviewed safe responses is necessary before this system is placed in front of anyone and is the first piece of work we recommend \cite{who2023suicide}. Next, a component-wise ablation where the three things the pipeline provides at once (emotion choice, per emotion guidance and the attention cues) are separated, and the contribution of each is known instead of inferred. The rule-adherence gap shall be bridged by checking and repairing the reply after it's written, and repetition shall be controlled at the decoding time, both of which were intentionally deferred here to keep the evaluated outputs fixed.

\textbf{Beyond English:} The decoupled structure is an advantage here, since a new language requires a new detector and a new guidance but not a new generator and the detector is only 125 million parameters. Three steps follow: collecting or adapting the emotional conversations in Bangla, including the code-mixed Bangla-English that people here actually type; retraining the detector from a multilingual encoder; and rewriting the per-emotion guidance with speakers of the language, as the nature of the question and answers would differ based on location and a translated prompt would bring English assumptions into the response. The literature of emotion classification with low resources indicates that the focus of the difficulty will be on the detector \cite{ed_github}.

\textbf{Stronger evaluation:} The evaluations should be made not only on individual responses but also on entire conversations evaluated by human evaluators and the human evaluation should be repeated at a size that allows statistical testing, ideally with mental-health professionals as evaluators. More members in the empathy panel would be nice, and the numbers of cost should be repeated on other hardware, such as machines without a graphics card, since it's about deployment.

\textbf{Final remarks:} This work is not about a new record on a benchmark and it wasn't designed to get one. It's in demonstrating, through measurement, not argument, that a system built from a small specialised model and a modest general one can be more emotionally responsive than a single model of the same size, can run on hardware already owned by the student so that a person's disclosure never leaves their machine, and can show exactly what it understood before it answers. How to arrange them will be more important than how big to make them, and a design whose emotional reasoning can be inspected and corrected and replaced one part at a time will be what a sensitive application will require, as capable open models continue to get smaller.
